%% gateways.tex
\section{Integrative Flux Bounds}\label{sec:gateways}

Conservation alone leaves the flux distribution underdetermined: the
mole-balance constraint fixes a subspace, not a point, and the box of
bounds together with the objective is what narrows that subspace to a
prediction. The objective is a modeling choice. The bounds are where
biology enters, and they
enter not as numbers to be looked up but as a model in their own right.
Writing a single upper bound as a product of factors makes that model
visible. For reaction $j$, the factorized bound is:
\begin{equation}\label{eq:general-bound}
  -\delta_j\Bigl[V^{\circ}_{\max,j}\,(e/e^{\circ})\,\theta_j(\cdot)\,f_j(\cdot)\Bigr]
  \;\le\; \hat v_j \;\le\;
  V^{\circ}_{\max,j}\,(e/e^{\circ})\,\theta_j(\cdot)\,f_j(\cdot),
\end{equation}
where $V^{\circ}_{\max,j}$ sets a reference capacity carrying the units
of flux and the remaining factors are dimensionless corrections that
scale that capacity up or down. Each factor represents a distinct kind
and level of biological information: $\delta_j$ encodes thermodynamic
directionality, $V^{\circ}_{\max,j}$ and $f_j$ encode enzyme kinetics,
$(e/e^{\circ})$ encodes gene expression, and $\theta_j$ encodes
allosteric regulation of enzyme activity. A bound therefore provides an
explicit accounting of what is known about a reaction. With no
information beyond a capacity scale the correction factors are set to
unity, and each is refined as measurements or mechanistic models become
available. This is also where the underdetermined polytope is finally
narrowed: not by conservation but by the bounds, and every informative
factor tightens that polytope using evidence of a particular kind, so that
tracing a prediction back to its cause is a matter of asking which
factor in Equation~\eqref{eq:general-bound} was responsible.

The thermodynamic factor is the switch $\delta_j$, which decides
whether the reaction may carry flux in the reverse direction at all. It
takes one of two values:
\begin{equation}\label{eq:thermo-gate}
  \delta_j \;=\;
  \begin{cases}
    1, & \text{reaction $j$ reversible},\\[2pt]
    0, & \text{reaction $j$ irreversible},
  \end{cases}
\end{equation}
so that $\delta_j=1$ leaves the lower bound of
Equation~\eqref{eq:general-bound} negative and the reaction free to run
backward, while $\delta_j=0$ collapses the lower bound to zero and pins
the reaction to the forward direction. The assignment follows the sign
of $\Delta G^{\circ}_j-\Delta G^{*}_j$, the standard free energy of
reaction measured against a threshold $\Delta G^{*}_j$ meant to
represent the physiological concentration range, or equivalently a
cutoff on the equilibrium constant $K_{\mathrm{eq},j}$: a reaction whose
driving force cannot plausibly change sign anywhere in that range is
declared irreversible, and its reverse flux is forbidden. This is a
heuristic directionality rule, a stand-in for propagating a stated
concentration range through $\Delta G=\Delta G^{\circ}+RT\ln Q$ and
checking directly whether the sign can flip, not a substitute for that
calculation where the range is actually known. Standard free energies
for this test are drawn from group- and component-contribution
estimates, for which eQuilibrator provides genome-scale coverage
\cite{Beber2022}.

When the reaction is allowed to proceed, its capacity is set by enzyme
kinetics, and this is the role of the reference scale
$V^{\circ}_{\max,j}$ together with the saturation factor $f_j$. The
scale is a maximal velocity, given by the product of a turnover number
and a characteristic enzyme abundance:
\begin{equation}\label{eq:vmax}
  V^{\circ}_{\max,j} \;=\; k^{\circ}_{\mathrm{cat},j}\,e^{\circ},
\end{equation}
where $k^{\circ}_{\mathrm{cat},j}$ is the catalytic rate constant of the
enzyme catalyzing reaction $j$ and $e^{\circ}$ is a reference enzyme
concentration. Turnover numbers are tabulated per enzyme and substrate
in BRENDA \cite{Chang2021}, and a representative $e^{\circ}$ can be read
from compilations of cellular abundances such as BioNumbers
\cite{Milo2010}. The scale $V^{\circ}_{\max,j}$ is the flux the enzyme
would sustain at reference abundance and saturating substrate, and the
fraction of that ceiling actually reached at the prevailing substrate
concentration is the saturation factor. That factor follows from the
enzyme mechanism itself. A single-substrate enzyme binds substrate
reversibly and turns it over to product according to the mechanism:
\begin{equation}\label{eq:mm-mechanism}
  E + S_j \;\underset{k_{-1}}{\overset{k_1}{\rightleftharpoons}}\; [ES]
  \;\overset{k_{\mathrm{cat}}}{\longrightarrow}\; E + P,
\end{equation}
with $[ES]$ the enzyme-substrate complex. Holding the complex at
quasi-steady state, on the standard argument that it equilibrates fast
relative to the depletion of substrate, and eliminating the free enzyme
through the conservation of total enzyme $[E]_T=[E]+[ES]$, gives the
complex as a hyperbolic function of substrate alone,
$[ES]=[E]_T\,[S_j]/(K_{M,j}+[S_j])$ with
$K_{M,j}\equiv(k_{-1}+k_{\mathrm{cat}})/k_1$.
The reaction velocity is the rate of the turnover step,
$v_j=k_{\mathrm{cat}}[ES]$, so writing the capacity as
$V_{\max,j}=k_{\mathrm{cat}}[E]_T$ recovers the Michaelis--Menten rate
law \cite{MichaelisMenten1913}, and the saturation factor is obtained by
normalizing the rate to that capacity:
\begin{equation}\label{eq:mm}
  v_j \;=\; V_{\max,j}\,\frac{[S_j]}{K_{M,j}+[S_j]}
  \qquad\Longrightarrow\qquad
  f_j([S_j]) \;=\; \frac{v_j}{V_{\max,j}}
  \;=\; \frac{[S_j]}{K_{M,j}+[S_j]}\;\in[0,1],
\end{equation}
with $[S_j]$ the substrate concentration and $K_{M,j}$ the Michaelis
constant. The saturation factor $f_j$ interpolates between negligible
capacity ($f_j\to0$ as substrate is exhausted) and the maximal capacity
($f_j\to1$ at saturation), and it is the same hyperbolic law that
governs single-enzyme assays. Because $f_j$
depends on the substrate concentration, incorporating this factor makes the
bound a function of the metabolic state rather than a fixed number, the
point at which a static flux balance begins to acquire the
state-dependence that a dynamic treatment carries to its conclusion.

The scale $V^{\circ}_{\max,j}$ was fixed at a reference enzyme abundance
$e^{\circ}$, but the abundance actually present is itself a variable set
by gene expression, and the ratio $(e/e^{\circ})$ is the factor through
which that variability enters. To treat abundance as a modeled
quantity rather than a fixed parameter, write balances for the
transcript $m_j$ and the protein $p_j$ of the gene encoding the enzyme:
\begin{align}
  \dot m_j &= r_{X,j}\,u_j - (\theta_{m,j}+\mu)\,m_j + \lambda_j,
  \label{eq:mrna}\\
  \dot p_j &= r_{L,j}\,w_j\,m_j - (\theta_{p,j}+\mu)\,p_j,
  \label{eq:protein}
\end{align}
where $r_{X,j}$ and $r_{L,j}$ are transcription and translation rate
capacities, $u_j$ and $w_j$ are dimensionless control functions that
modulate those capacities, and $\theta_{m,j}$ and $\theta_{p,j}$ are
first-order degradation rate constants for message and protein, and
$\lambda_j$ is a basal, or leak, source of transcript. The loss term in
each balance is not the degradation constant alone but the sum
$(\theta_{m,j}+\mu)$ or $(\theta_{p,j}+\mu)$: every intracellular
species is removed both by its own turnover and by dilution into the
growing culture at the specific growth rate $\mu$. That $\mu$ is not a
new parameter. It is the very growth-dilution term retained in the
constraint retaining growth dilution,
Equation~\eqref{eq:dilution-constraint}, carried
through consistently: an
enzyme is diluted by growth exactly as a metabolite is, and the same
$\mu$ that balances metabolic production in the flux constraint sets
the residence time of the transcripts and proteins that fix the flux
bounds. Holding the balances at steady state gives the closed-form
abundances:
\begin{equation}\label{eq:expression-ss}
  m^{\star}_j = \frac{r_{X,j}u_j+\lambda_j}{\theta_{m,j}+\mu},
  \qquad
  p^{\star}_j = \frac{r_{L,j}w_j\,m^{\star}_j}{\theta_{p,j}+\mu},
  \qquad
  (e/e^{\circ}) = \frac{p^{\star}_j}{e^{\circ}},
\end{equation}
and the protein level $p^{\star}_j$ is mapped onto the abundance of the
catalytic activity for reaction $j$ through the gene-protein-reaction
(GPR) rules of the reconstruction, which encode whether an activity
requires a single gene product, several assembled into a complex, or
any one of several interchangeable isozymes. Through
Equation~\eqref{eq:expression-ss} the expression factor couples the
flux ceiling to the entire upstream apparatus of transcription and
translation, with $\mu$ standing in every denominator as the visible
signature of the growth-aware bookkeeping begun in the mole-balance
derivation. The ratio $(e/e^{\circ})$ is also the
natural port for measured data: a transcriptomic or proteomic profile
enters here, rescaling the capacity of each reaction by the abundance
actually observed under a given condition and turning one
reconstruction into a family of condition-specific models.

One statistical mechanical construction can represent allosteric
activity, transcription, and translation. The common idea is to
classify configurations by function. Some configurations of an enzyme,
promoter, or translation-initiation complex permit the process of
interest; the remaining configurations do not. The physical
configurations and the criterion for activity differ among the three
processes, but the calculation is the same: the numerator sums the
configurations that permit the process, while the denominator sums all
configurations available to the regulated element.

Let $\mathcal S_q$ denote the set of configurations of a regulated
element associated with process $q$, where $q$ may be catalysis,
transcription, or translation. Configuration $s\in\mathcal S_q$ has
energy $\epsilon_s$, with its statistical weight given by:
\begin{equation}\label{eq:partition}
  \omega_s(\mathbf x)
    = g_s(\mathbf x)e^{-\beta\epsilon_s},
  \qquad
  Z_q(\mathbf x)
    = \sum_{r\in\mathcal S_q}\omega_r(\mathbf x),
  \qquad
  p_s(\mathbf x)
    = \frac{\omega_s(\mathbf x)}{Z_q(\mathbf x)},
\end{equation}
where $\beta=1/(k_BT)$ and $Z_q$ is the partition function. The factor
$g_s(\mathbf x)$ accounts for ligand concentrations, binding
stoichiometry, and degeneracy. For example, a configuration containing
$n$ bound molecules of an effector $x$ contributes a concentration
factor proportional to $(x/K)^n$. The partition function is therefore a
concentration-dependent sum over every configuration, not only those
that are functional.

Now let $\mathcal A_q\subseteq\mathcal S_q$ be the configurations that
permit process $q$, and let $\mathcal I_q=\mathcal S_q\setminus\mathcal
A_q$ be the configurations that do not. The dimensionless control for
that process is the probability of occupying the active subset:
\begin{equation}\label{eq:active-state-control}
  \Phi_q(\mathbf x)
    = \sum_{s\in\mathcal A_q}p_s(\mathbf x)
    = \frac{\displaystyle\sum_{s\in\mathcal A_q}
        g_s(\mathbf x)e^{-\beta\epsilon_s}}
      {\displaystyle\sum_{r\in\mathcal S_q}
        g_r(\mathbf x)e^{-\beta\epsilon_r}}
    = \frac{Z_{\mathcal A_q}}
      {Z_{\mathcal A_q}+Z_{\mathcal I_q}}.
\end{equation}
Here $Z_{\mathcal B}=\sum_{s\in\mathcal B}\omega_s$ denotes a
partition sum restricted to the subset $\mathcal B$. This ratio makes
the functional classification explicit. Every active
configuration appears in both the numerator and denominator; inactive
configurations appear only in the denominator. Consequently
$0\leq\Phi_q\leq1$, and an effector increases or decreases the control
according to whether it preferentially weights active or inactive
configurations.

The compact constants used in promoter models can be connected directly
to these Boltzmann weights. Choose configuration $0$ as a reference and
write $\Delta\epsilon_s=\epsilon_s-\epsilon_0$. The relative energetic
weight and the complete state weight are then related by:
\begin{equation}\label{eq:effective-state-weight}
  K_{\star,s}
    \equiv \frac{\omega_s/g_s}{\omega_0/g_0}
    = e^{-\beta\Delta\epsilon_s},
  \qquad
  \frac{\omega_s}{\omega_0}
    = \frac{g_s(\mathbf x)}{g_0(\mathbf x)}K_{\star,s}.
\end{equation}
With the conventional choice $\omega_0=g_0=1$, this becomes
$\omega_s=g_s(\mathbf x)K_{\star,s}$. Thus, a term such as $K_2f_{TL}$
combines an energetic factor $K_2=e^{-\beta\Delta\epsilon_2}$ with the
availability $f_{TL}$ of the ligand-bound transcription factor. In a
fitted promoter model, however, $K_1$ and $K_2$ are effective,
dimensionless weights: they may absorb a constant RNA-polymerase
concentration, standard-state factors, and other lumped contributions,
and need not equal measured dissociation constants.

Figure~\ref{fig:active-subset} summarizes the construction from molecular
configurations to process-specific control:
\begin{figure}[H]
  \centering
  \includegraphics[width=0.88\textwidth]{active_subset.pdf}
  \caption{Active-subset construction of biological control factors.
  (a) Configuration weights combine relative energetics with regulator
  availability; blue states permit the process and gray states do not.
  (b) Active states enter the numerator, while all states enter the
  partition-function denominator. The promoter shown is schematic.
  (c) Specifying the relevant configurations gives allosteric activity
  $\theta_j$, transcriptional control $u_j$, or translational control
  $w_j$ from the same statistical accounting.}
  \label{fig:active-subset}
\end{figure}

The inducible-promoter models of Moon, Voigt, and coworkers provide a
concrete example of this construction \cite{Moon2012}. Their
Supplementary Figure~S7 and Equations~3--6 enumerate promoter
configurations and assign a partition-function term to each one. For
the Lux promoter, the configurations have weights
$1$, $K_1$, and $K_2f_{TL}$, corresponding respectively to an empty
promoter, RNA-polymerase binding, and joint binding of RNA polymerase
with the ligand-bound LuxR complex. The two RNA-polymerase-bound
configurations permit transcription and therefore appear in the
numerator. For the Tac promoter, the empty, RNA-polymerase-bound, and
repressor-bound configurations have weights $1$, $K_1$, and $K_2f_T$;
only the RNA-polymerase-bound configuration permits transcription.
After normalization by the maximal promoter activity, the two controls
are therefore given by:
\begin{equation}\label{eq:moon-promoter-control}
  u_{\mathrm{Lux}}
    = \frac{K_1+K_2f_{TL}}{1+K_1+K_2f_{TL}},
  \qquad
  u_{\mathrm{Tac}}
    = \frac{K_1}{1+K_1+K_2f_T}.
\end{equation}
Here $f_{TL}$ is the ligand-bound fraction of the transcription factor,
whereas $f_T$ is its ligand-free fraction. The Lux example adds weight
to an active configuration as inducer increases; the Tac example
removes weight from an inactive, repressor-bound configuration. Both
are direct specializations of Equation~\eqref{eq:active-state-control}.

A two-configuration element illustrates the reduction. If an inactive
configuration has weight $\omega_I=1$ and an effector-stabilized active
configuration has weight
$\omega_A=(x/K)^n e^{-\beta\Delta\epsilon}$, then the control becomes:
\begin{equation}\label{eq:two-state-control}
  \Phi_q(x)
    = \frac{(x/K)^n e^{-\beta\Delta\epsilon}}
      {1+(x/K)^n e^{-\beta\Delta\epsilon}},
\end{equation}
where $\Delta\epsilon$ is the active-state energy relative to the
inactive reference. The familiar Hill-like form is thus a normalized
active-state weight.
If the effector instead stabilizes an inactive configuration, its
concentration-dependent weight enters the denominator but not the
numerator, and the control decreases with $x$.

The process-specific factors used above are instances of
Equation~\eqref{eq:active-state-control}:
\begin{equation}\label{eq:control-mapping}
  \theta_j = \Phi_{\mathrm{cat},j},
  \qquad
  \bar u_j = \Phi_{\mathrm{TX},j},
  \qquad
  w_j = \Phi_{\mathrm{TL},j}.
\end{equation}
The allosteric factor $\theta_j$ multiplies catalytic capacity directly
in Equation~\eqref{eq:general-bound}. The total transcriptional control
$\bar u_j$ is separated below into regulated and basal contributions,
while $w_j$ multiplies translational capacity in
Equation~\eqref{eq:protein}. These dimensionless controls are distinct
from the degradation-rate constants $\theta_{m,j}$ and
$\theta_{p,j}$.

Basal expression follows by subdividing the transcriptionally active
set. Let $\mathcal A_{\mathrm{TX}}^{\mathrm{reg}}$ contain active
configurations reached through the regulated input and let
$\mathcal A_{\mathrm{TX}}^{\mathrm{basal}}$ contain configurations that
support transcription without that input. Their contributions are:
\begin{equation}\label{eq:control-split}
  \bar u_j
    = \frac{Z_{\mathcal A_{\mathrm{TX}}^{\mathrm{reg}}}}{Z_{\mathrm{TX}}}
    + \frac{Z_{\mathcal A_{\mathrm{TX}}^{\mathrm{basal}}}}{Z_{\mathrm{TX}}}
    = u_j + u_j^{\dagger},
  \qquad
  \lambda_j \equiv r_{X,j}u_j^{\dagger}.
\end{equation}
The regulated contribution $u_j$ enters the transcription term in
Equation~\eqref{eq:mrna}, whereas the basal contribution
$u_j^{\dagger}$ supplies its leak source $\lambda_j$. The promoter-state
specialization follows the construction used by Moon, Voigt, and
coworkers \cite{Moon2012}. The effective biophysical model of Adhikari and
coworkers extends the same active-subset accounting to transcriptional
and translational control \cite{Adhikari2020}; the factorization here
places those controls alongside the corresponding allosteric factor.

Each information layer requires data that is often unavailable:
thermodynamic estimates, turnover numbers, substrate concentrations,
expression measurements, and regulatory parameters. When that data is
scarce the correction factors are set to unity,
$f_j\approx1$, $(e/e^{\circ})\approx1$, and
$\theta_j\approx1$, and the general bound collapses to the two-parameter
box:
\begin{equation}\label{eq:simplified-bound}
  -\delta_j\,V^{\circ}_{\max,j} \;\le\; \hat v_j \;\le\; V^{\circ}_{\max,j}.
\end{equation}
The simplified form in Equation~\eqref{eq:simplified-bound} retains only
the thermodynamic factor and the capacity scale, and it is the reversible or irreversible
$\pm V_{\max}$ box used throughout genome-scale flux balance analysis.
It is the tractable baseline representation of a bound when
nothing beyond a capacity estimate is known. The factorized form in
Equation~\eqref{eq:general-bound} records which
assumptions that baseline makes and which factor requires additional
information to relax each of them.

Each class of information just described has a named, established predecessor in the
constraint-based modeling literature, and this chapter's contribution is
the unifying notation and workflow that connects them, not the invention
of any individual constraint. The thermodynamic factor $\delta_j$ is the
central device of thermodynamics-based flux analysis, which propagates
standard free energies and metabolite activity ranges through the whole
network to rule out infeasible flux and concentration profiles together
\cite{Henry2007}. The kinetic factors, $V^{\circ}_{\max,j}$ built from a
turnover number and an enzyme abundance together with the saturation
factor $f_j$, are the organizing idea behind enzyme-constrained flux
balance analysis, most fully realized in the GECKO framework, which
folds a proteome-wide capacity constraint into a genome-scale
reconstruction \cite{Sanchez2017}. The expression factor $(e/e^{\circ})$,
transcript and protein balances entered as a modeled quantity rather
than a fixed parameter, is the metabolism-and-expression (ME) model
construction, which couples the transcription and translation machinery
directly to the metabolic network at genome scale
\cite{OBrien2013}. The expression control functions $u_j$ and $w_j$
connect to two established lines of regulatory constraint-based
modeling: regulatory flux balance analysis, which switches expression-linked
reactions on and off with Boolean rules derived from a transcriptional
regulatory network \cite{Covert2001}, and PROM, which replaces the
Boolean rule with a continuous, data-derived probability of an
interaction being active \cite{Chandrasekaran2010}. The allosteric
factor $\theta_j$ represents the distinct, direct regulation of enzyme
activity after the enzyme has been expressed. What the integrative
factorization of Equation~\eqref{eq:general-bound} adds is not a fifth
constraint class but a common accounting of how these four
already-established devices contribute to one bound, so that a single
notation and a single worked derivation carry across all four rather
than requiring a separate formalism, and a separate paper, for each.

The expression model above lumped transcription and
translation into aggregate capacities $r_{X,j}$ and $r_{L,j}$, but these
processes can themselves be resolved to the level of the sequences they
read. Transcription of a gene and translation of its protein consume
nucleotides, amino acids, and energy in stoichiometric ratios fixed by
the gene and protein sequences, and those precursor demands can be
computed directly from the sequences and entered into $\mathbf S$ as
reactions in their own right \cite{AllenPalsson2003}, with the
precursors supplied by exchange reactions across the system boundary.
Expression then becomes flux-bearing chemistry inside the same linear
program as metabolism, a construction carried out for cell-free protein
synthesis by Vilkhovoy and coworkers \cite{Vilkhovoy2018}, and it is
the mechanism behind the cell-free capstone taken up later in this
chapter.
